\section{Structural Emergence of Mean Bias}

We first establish the global empirical pattern of \emph{mean bias} and then explain its mechanism. Here, mean bias refers to a coherent, nonzero feature-wise mean component that increasingly dominates activation energy across layers and training stages. This motivates a three-stage causal account: frequency-weighted initialization, nonlinear regeneration/amplification, and residual accumulation.

\paragraph{Empirical Trajectory: Mean-Bias Energy Rises with Training and Depth.}

We quantify mean magnitude by the normalized ratio
\begin{equation}
R = \frac{\|\mu\|_2}{\sqrt{\mathbb{E}\|X\|^2/N}}.
\end{equation}

Using intermediate activation matrices from Qwen3-0.6B at early (10k) and late (170k) checkpoints, we decompose activations into mean + spike + tail components and track mean-energy share across embedding, shallow, middle, and deep layers. Figure~\ref{fig:energy_grid} shows a clear trend: shallow layers already exhibit strong mean dominance early, while deep layers become increasingly mean-dominated later; in all profiled layers, mean-bias energy increases over training.

\begin{figure*}[h]
\centering
\begin{subfigure}{0.49\textwidth}
\includegraphics[width=\textwidth]{figures/energy/energy_10k.png}
\caption{Early stage (10k)}
\end{subfigure}\hfill
\begin{subfigure}{0.49\textwidth}
\includegraphics[width=\textwidth]{figures/energy/energy_170k.png}
\caption{Late stage (170k)}
\end{subfigure}
\caption{Mean-bias energy decomposition across layers and training stages, shown as compact summaries for early (10k) and late (170k) checkpoints.}
\label{fig:energy_grid}
\end{figure*}

\paragraph{Mechanism: A Three-Stage Causal Chain.}

\noindent \emph{Stage I: Statistical Origin via Frequency-Weighted Embedding Expectation.}
A seed bias appears at the input representation level. Let $\mathcal{V}$ denote vocabulary and $p(v)$ token frequency. The corpus-level expected embedding is
\begin{equation}
\mu_{\text{embed}} = \mathbb{E}_{v\sim p}[E_v] = \sum_{v\in\mathcal{V}} p(v)E_v,
\label{eq:embedding_mean}
\end{equation}
where $E_v\in\mathbb{R}^H$ is the embedding of token $v$ \cite{mikolov2013distributed,arora2017simple}. Because $p(v)$ is highly non-uniform (Zipf-like), high-frequency tokens receive disproportionate updates and induce coherent embedding alignment, forming a nontrivial mean component before deeper processing.

\noindent \emph{Stage II: Nonlinear Regeneration and Amplification in Attention/FFN.}
Even when pre-activation features are centered, non-odd nonlinearities regenerate a nonzero mean. For a symmetric zero-mean random variable $z$, non-odd activations satisfy
\begin{equation}
\mathbb{E}[\phi(z)] > 0,
\end{equation}
which holds exactly for ReLU and approximately for GELU/SwiGLU due to asymmetric suppression of negative responses. In parallel, softmax attention favors coherent high-similarity directions through exponential reweighting, increasing mass on already aligned components. Empirically, operator-level input/output comparisons (Figure~\ref{fig:operator_change}) show a consistent mean-energy increase after nonlinear processing: at 10k steps, attention softmax strengthens mean dominance; in the FFN path, SwiGLU further amplifies this effect.

\begin{figure}[t]
\centering
\begin{subfigure}{0.49\textwidth}
\includegraphics[width=\textwidth]{figures/energy/layer0_attn_input_output_energy_10k_notitle.png}
\caption{Softmax attention (10k)}
\end{subfigure}\hfill
\begin{subfigure}{0.49\textwidth}
\includegraphics[width=\textwidth]{figures/energy/layer0_ffn_input_output_energy_combo.png}
\caption{FFN SwiGLU}
\end{subfigure}
\caption{Operator-level input/output energy comparison. Left: attention softmax at 10k. Right: FFN SwiGLU. Both operators increase mean-bias energy from input to output.}
\label{fig:operator_change}
\end{figure}

\noindent \emph{Stage III: Residual Preservation and Network-Wide Accumulation.}
Residual connections preserve and propagate coherent mean components across depth:
\begin{equation}
\mu_{l+1}=\mu_l+\Delta\mu_l,
\label{eq:residual_accumulation}
\end{equation}
where $\Delta\mu_l$ is the layer-induced mean shift. This additive pathway prevents cancellation of existing bias and enables cumulative growth through the network. In high dimensions, even modest coordinate-wise drift becomes large in norm:
\begin{equation}
\|\mu\|_2 \sim \sqrt{H}\,\bar{\mu},
\label{eq:mean_scaling}
\end{equation}
with hidden size $H$ \cite{tropp2012user,halko2011randomized}. This geometric scaling explains why mean-dominated components contribute disproportionately to extreme activations and quantization-sensitive outliers.
